%%%%%
%
\section{Passage du rep\`ere ECEF au rep\`ere ENU}

\itemstyle{Op\'erateur~:} \operfmt{ecef\_to\_enu}
\index{ecef\_to\_enu}

%%%%%OBJET%%%%
\itemstyle{Objet~:}
Transformation d'un vecteur du rep\`ere ECEF
au rep\`ere ENU dans un cadre simpliste (terre sph\'erique).

%%%%%ARGUMENTS%%%%
\bigskip
%%%%%ENTREES%%%%
\begin{tabarg}
\multicolumn{3}{l}{\titretab {Entr\'ees}}\\
\hline
$\lambda$ & Vecteur, {\tt float} & Latitude (\textit{degr\'es})\\
$\phi$ & Vecteur, {\tt float} & Longitude (\textit{degr\'es})\\
$V_x$ & Vecteur, {\tt float} & Composante suivant l'axe X du rep\`ere ECEF\\
$V_y$ & Vecteur, {\tt float} & Composante suivant l'axe Y du rep\`ere ECEF\\
$V_z$ & Vecteur, {\tt float} & Composante suivant l'axe Z du rep\`ere ECEF\\
\end{tabarg}

\medskip
%%%%%SORTIES%%%%
\begin{tabarg}
\multicolumn{3}{l}{\titretab {Sorties}}\\
\hline
$V_e$ & Vecteur, {\tt float} & Composante Ouest-Est (m\^eme unit\'e que $V_x$)\\
$V_n$ & Vecteur, {\tt float} & Composante Sud-Nord (m\^eme unit\'e que $V_x$)\\
$V_v$ & Vecteur, {\tt float} & Composante Verticale (m\^eme unit\'e que $V_x$)\\
\end{tabarg}


%%%%%FORMULE%%%%
\itemstyle{Formule}

%%ecef_to_enu
\begin{eqnarray*}
V_v & = & cos \lambda . ( cos \phi.V_x + sin \phi V_y + tan \lambda . V_z) \\
V_e & = & cos \phi . V_y - sin \phi .V_x \\
V_n & = & \frac{V_z}{cos \lambda} - tan \lambda . V_v \\
\end{eqnarray*}


%%%% COMMENTAIRE %%%%

\itemstyle{Commentaire~:}

Le rep\`ere ECEF (Earth Centered Earth Fixed) est d\'efini
de la mani\`ere suivante~:
\begin{description}\itemsep0pt
 \item[~ Origine :] Centre de la Terre
 \item[~ X :] de l'origine vers le point $0^\circ$ N $0^\circ$ W
 \item[~ Y :] de l'origine vers le point $0^\circ$ N $90^\circ$ E
 \item[~ Z :] de l'origine vers le p\^ole Nord 
\end{description}

Le rep\`ere ENU (East, North, Up : rep\`ere "g\'eographique" ) est d\'efini
de la mani\`ere suivante~:
\begin{description}\itemsep0pt
 \item[~ Origine :] Position du mobile
 \item[~ X :] de l'origine vers l'Est tangent \`a la surface terrestre
 \item[~ Y :] de l'origine vers le Nord tangent \`a la surface terrestre
 \item[~ Z :] de l'origine vers le haut perpendiculaire \`a la surface terrestre
\end{description}

Le cas limite (les p\^oles, o\`u le rep\`ere ENU n'est pas d\'efini)
n'est pas trait\'e.

%%%%%AUTEUR%%%%
\itemstyle{Auteur~:}
CNRM/GMEI/TRAMM

%%%%%REFERENCE%%%%
\itemstyle{R\'ef\'erences~:}
Hofmann-Wellenhof \cite{Hofmann-GPS}
